\section{Rabin--Karp Algorithm}
\label{sec:rabin-karp}

All the algorithms we have seen so far---naive, KMP, Boyer--Moore,
Aho--Corasick---are \emph{comparison-based}: their main primitive is
comparing two characters.
The Rabin--Karp algorithm takes a completely different approach: it
replaces character comparisons with \emph{arithmetic operations}.
The idea is simple and elegant: interpret strings as numbers, and check
whether the ``number'' of the pattern equals the ``number'' of the
current window in the text.

\aside{Rabin--Karp was published by Richard Karp and Michael Rabin in
1987, although the ideas had been circulating since the late 1970s.
It is historically important as one of the first \emph{randomised}
algorithms for string matching.}

\subsection{Strings as Numbers}

Without loss of generality, we can assume $\Sigma = \{0, 1, \ldots,
d-1\}$ where $d = |\Sigma|$.
If the original alphabet is not numeric, we simply map each character
to a digit: for instance, with $\Sigma = \{\texttt{a}, \texttt{b},
\texttt{c}, \texttt{d}\}$ we set $\texttt{a} \mapsto 0$,
$\texttt{b} \mapsto 1$, $\texttt{c} \mapsto 2$, $\texttt{d} \mapsto 3$.
This mapping takes $\Oh{n + m}$ time and is done once as a
preprocessing step.

% ── Figure: string-to-number mapping ─────────────────────
\begin{figure}[H]
\centering
\begin{tikzpicture}[
    cell/.style={minimum width=8mm, minimum height=8mm,
                 font=\ttfamily\small, inner sep=0pt, text centered,
                 draw=gray!40, rounded corners=1.5pt},
    C/.style={cell, fill=blue!6},
    N/.style={cell, fill=teal!10, text=teal!70!black},
]
\def\cw{1.05}

% ── S row ─────────────────────────────────────────────────
\node[font=\small\sffamily\bfseries, text=gray!60, anchor=east]
  at ({-0.4*\cw}, 0) {$S =$};
\foreach \i/\c in {1/b, 2/a, 3/d, 4/c, 5/c, 6/a, 7/b, 8/c, 9/d, 10/b}
  \node[C] (S\i) at ({\i*\cw}, 0) {\c};

% ── arrow ──────────────────────────────────────────────────
\draw[-Stealth, thick, gray!60] ({5.5*\cw}, -0.65) -- ({5.5*\cw}, -1.35);

% ── S' row ─────────────────────────────────────────────────
\node[font=\small\sffamily\bfseries, text=gray!60, anchor=east]
  at ({-0.4*\cw}, -2) {$S' =$};
\foreach \i/\c in {1/1, 2/0, 3/3, 4/2, 5/2, 6/0, 7/1, 8/2, 9/3, 10/1}
  \node[N] (N\i) at ({\i*\cw}, -2) {\c};

% ── alphabet labels ──────────────────────────────────────
\node[font=\small\sffamily, text=gray!50, anchor=west]
  at ({11.2*\cw}, 0) {$\Sigma = \{\texttt{a,b,c,d}\}$};
\node[font=\small\sffamily, text=gray!50, anchor=west]
  at ({11.2*\cw}, -2) {$\Sigma' = \{0,1,2,3\}$};
\end{tikzpicture}
\caption{%
  Mapping a string $S$ over $\Sigma = \{\texttt{a,b,c,d}\}$ to a
  numeric string $S'$ over $\Sigma' = \{0,1,2,3\}$.
  The mapping is $\texttt{a}\mapsto 0$, $\texttt{b}\mapsto 1$,
  $\texttt{c}\mapsto 2$, $\texttt{d}\mapsto 3$.
}
\label{fig:rk-char-to-digit}
\end{figure}

Once both $T$ and $P$ are expressed as sequences of digits in base $d$,
we can interpret any length-$m$ window of $T$ as a number.
Formally, define for each position $s \in \{1, \ldots, n - m + 1\}$:
\[
  t_s \;=\; \bigl(T[s]\cdot d^{m-1} \;+\; T[s+1]\cdot d^{m-2}
           \;+\; \cdots \;+\; T[s+m-1]\cdot d^{0}\bigr)
\]
and similarly for the pattern:
\[
  p \;=\; \bigl(P[1]\cdot d^{m-1} \;+\; P[2]\cdot d^{m-2}
          \;+\; \cdots \;+\; P[m]\cdot d^{0}\bigr).
\]
Clearly, there is an occurrence of $P$ at position $s$ if and only if
$t_s = p$.
If we could compare these numbers in $\Oh{1}$ time, one scan through
$T$ would give us an $\Oh{n}$ algorithm.

\aside{The key obstacle is that $t_s$ and $p$ can be
\emph{astronomically large}: a string of $m$ characters in base~$d$
represents a number up to $d^m - 1$, which requires $\Om{m \log d}$
bits.
On a standard RAM machine, comparing two numbers of $m$ digits takes
$\Oh{m}$ time---no better than character-by-character comparison.}

\subsection{The Fingerprinting Idea}

The Rabin--Karp approach side-steps the large-number problem by working
\emph{modulo a prime} $q$.
Instead of comparing $t_s$ and $p$ directly, we compare their
\textbf{fingerprints}:
\[
  t_s \bmod q
  \quad\stackrel{?}{=}\quad
  p \bmod q.
\]
Since the fingerprints are numbers in $\{0, 1, \ldots, q-1\}$, each
comparison takes $\Oh{1}$ time (assuming $q$ fits in a machine word).

Two things can happen:

\begin{enumerate}[nosep]
  \item \textbf{$t_s \bmod q \neq p \bmod q$}: then certainly
        $t_s \neq p$, so there is \textbf{no match} at position~$s$.
        This is a \emph{true negative}---we can safely skip this
        position.

  \item \textbf{$t_s \bmod q = p \bmod q$}: this is a \textbf{hit}.
        Unfortunately, it does \emph{not} guarantee $t_s = p$:
        two different numbers can have the same remainder modulo $q$.
        We call this a \textbf{spurious hit} (or \emph{false positive}).
        When a hit occurs, we must verify character by character whether
        $T[s, s+m-1] = P$, which costs $\Oh{m}$.
\end{enumerate}

\aside{Think of it as a cheap ``quick reject'' filter: most positions
fail the modular test and are discarded instantly.
Only the rare hits require the expensive $\Oh{m}$ verification.}

\subsection{Rolling Hash via Horner's Rule}

Computing $p \bmod q$ from scratch is easy: we apply Horner's rule.
\[
  p \bmod q \;=\;
  \bigl(\!\bigl((P[1]\cdot d + P[2])\cdot d + P[3]\bigr)
  \cdot d + \cdots + P[m]\bigr) \bmod q.
\]
This takes $\Oh{m}$ multiplications and additions, all performed modulo
$q$ so that intermediate values never exceed $q \cdot d$.
The same computation gives us $t_1 \bmod q$ (the fingerprint of the
first window).

The crucial observation is that we can compute $t_{s+1} \bmod q$ from
$t_s \bmod q$ in $\Oh{1}$ time, without scanning $m$ characters again.
The relationship between consecutive windows is:
\begin{align*}
  t_{s+1}
  &= \bigl(T[s+1]\cdot d^{m-1} + T[s+2]\cdot d^{m-2} + \cdots
     + T[s+m]\cdot d^{0}\bigr) \\
  &= \bigl(t_s - T[s]\cdot d^{m-1}\bigr)\cdot d \;+\; T[s+m].
\end{align*}
In words: \emph{remove} the contribution of the leftmost character
$T[s]$, \emph{shift} everything one position to the left (multiply by
$d$), and \emph{add} the new rightmost character $T[s+m]$.

Taking everything modulo $q$:
\[
  t_{s+1} \bmod q
  \;=\;
  \Bigl[\bigl(t_s - T[s]\cdot h\bigr)\cdot d \;+\; T[s+m]\Bigr] \bmod q,
  \qquad
  \text{where } h = d^{m-1} \bmod q.
\]
The value $h = d^{m-1} \bmod q$ is precomputed once in $\Oh{m}$ time (or
$\Oh{\log m}$ using fast exponentiation).

% ── Figure: rolling hash ─────────────────────────────────
\begin{figure}[H]
\centering
\begin{tikzpicture}[
    cell/.style={minimum width=8mm, minimum height=8mm,
                 font=\ttfamily\small, inner sep=0pt, text centered,
                 rounded corners=1.5pt},
    C/.style={cell, draw=gray!40, fill=gray!6},
    W/.style={cell, draw=teal!50, fill=teal!10, text=teal!70!black},
    O/.style={cell, draw=red!50, fill=red!8, text=red!60!black},
    N/.style={cell, draw=violet!50, fill=violet!10, text=violet!60!black},
]
\def\cw{0.95}

% ── T row ─────────────────────────────────────────────────
\node[font=\small\sffamily\bfseries, text=gray!60, anchor=east]
  at ({-0.4*\cw}, 0) {$T$};
\foreach \i/\c in {1/3, 2/1, 3/4, 4/1, 5/5, 6/9, 7/2, 8/6}
  \node[C] (T\i) at ({\i*\cw}, 0) {\c};
% indices
\foreach \i in {1,...,8}
  \node[font=\tiny\sffamily, text=gray!50, above=2pt] at ({\i*\cw}, 0.42) {\i};

% ── Window s ──────────────────────────────────────────────
\node[font=\scriptsize\sffamily, text=gray, anchor=east] at ({-0.55*\cw}, -1.4) {$t_s$:};
\node[O] at ({2*\cw}, -1.4) {1};
\node[W] at ({3*\cw}, -1.4) {4};
\node[W] at ({4*\cw}, -1.4) {1};
\node[W] at ({5*\cw}, -1.4) {5};
% brace
\draw[teal!50!black, thick, decorate,
      decoration={brace, mirror, raise=6pt, amplitude=3pt}]
  ({1.55*\cw}, -1.85) -- ({5.45*\cw}, -1.85)
  node[midway, below=11pt, font=\scriptsize\sffamily, text=teal!50!black]
  {window at position $s$};

% ── Window s+1 ────────────────────────────────────────────
\node[font=\scriptsize\sffamily, text=gray, anchor=east] at ({-0.55*\cw}, -3.0) {$t_{s+1}$:};
\node[W] at ({3*\cw}, -3.0) {4};
\node[W] at ({4*\cw}, -3.0) {1};
\node[W] at ({5*\cw}, -3.0) {5};
\node[N] at ({6*\cw}, -3.0) {9};
% brace
\draw[violet!50!black, thick, decorate,
      decoration={brace, mirror, raise=6pt, amplitude=3pt}]
  ({2.55*\cw}, -3.45) -- ({6.45*\cw}, -3.45)
  node[midway, below=11pt, font=\scriptsize\sffamily, text=violet!50!black]
  {window at position $s+1$};

% ── Annotations ──────────────────────────────────────────
\draw[-Stealth, red!60!black, thick]
  ({2*\cw}, -0.5) -- ({2*\cw}, -1.0)
  node[midway, right=4pt, font=\scriptsize\sffamily, text=red!60!black]
  {remove};
\draw[-Stealth, violet!60!black, thick]
  ({6*\cw}, -0.5) -- ({6*\cw}, -2.6)
  node[midway, right=4pt, font=\scriptsize\sffamily, text=violet!60!black]
  {add};

\end{tikzpicture}
\caption{%
  Rolling hash: computing $t_{s+1}$ from $t_s$.
  The \textcolor{red!60!black}{\textbf{red}} digit $T[s]$ is
  \emph{removed} from the high end, the remaining digits are shifted
  left (multiplied by $d$), and the
  \textcolor{violet!60!black}{\textbf{violet}} digit $T[s+m]$ is
  \emph{added} at the low end.
  Each update costs $\Oh{1}$ arithmetic operations modulo $q$.
}
\label{fig:rk-rolling-hash}
\end{figure}

This ``rolling'' update is the heart of Rabin--Karp: after the initial
$\Oh{m}$ setup, each new window costs only $\Oh{1}$ to fingerprint.

\subsection{The Algorithm}

Putting the pieces together, the Rabin--Karp algorithm works in three
phases:

\begin{enumerate}[nosep]
  \item \textbf{Preprocessing.}
        Compute $h = d^{m-1} \bmod q$.
        Compute $p = P \bmod q$ using Horner's rule.
        Compute $t_1 = T[1,m] \bmod q$ using Horner's rule.
        This takes $\Oh{m}$ time.

  \item \textbf{Scanning.}
        For each position $s = 1, 2, \ldots, n-m+1$:
        \begin{itemize}[nosep]
          \item If $t_s \equiv p \pmod{q}$: \emph{hit} --- verify
                character by character whether $T[s,\,s+m-1] = P$.
                If yes, report occurrence at position~$s$.
          \item If $t_s \not\equiv p \pmod{q}$: skip (no match possible).
          \item Compute $t_{s+1}$ from $t_s$ using the rolling hash
                formula.
        \end{itemize}

  \item \textbf{Output.} All reported positions are true occurrences.
\end{enumerate}

\begin{algorithm}[H]
  \caption{Rabin--Karp}
  \label{alg:rabin-karp}
  \KwIn{Text $T[1\ldots n]$, Pattern $P[1\ldots m]$, alphabet size $d$,
        prime $q$}
  \KwOut{All positions $s$ where $P$ occurs in $T$}
  $h \gets d^{m-1} \bmod q$\;
  \BlankLine
  \tcp{Horner's rule for pattern fingerprint}
  $p \gets 0$\;
  \For{$j \gets 1$ \KwTo $m$}{
    $p \gets (p \cdot d + P[j]) \bmod q$\;
  }
  \BlankLine
  \tcp{Horner's rule for first window}
  $t \gets 0$\;
  \For{$j \gets 1$ \KwTo $m$}{
    $t \gets (t \cdot d + T[j]) \bmod q$\;
  }
  \BlankLine
  \tcp{Slide the window}
  \For{$s \gets 1$ \KwTo $n - m + 1$}{
    \If{$p = t$}{
      \tcp{Hit: verify character by character}
      $k \gets 0$\;
      \While{$k < m$ \textbf{and} $P[k+1] = T[s+k]$}{
        $k \gets k + 1$\;
      }
      \If{$k = m$}{
        \textbf{output} $s$\;
      }
    }
    \If{$s < n - m + 1$}{
      $t \gets \bigl((t - T[s] \cdot h) \cdot d + T[s+m]\bigr)
              \bmod q$\;
    }
  }
\end{algorithm}

\subsection{Correctness}

The correctness relies on two observations:
\begin{enumerate}[nosep]
  \item \textbf{No false negatives.}
        If $t_s \bmod q \neq p \bmod q$, then certainly $t_s \neq p$,
        so skipping position~$s$ is safe.
        We never miss a true occurrence.

  \item \textbf{Hits are verified.}
        When $t_s \bmod q = p \bmod q$, the algorithm performs a full
        character-by-character comparison.
        A position~$s$ is reported as an occurrence \emph{only if}
        $T[s, s+m-1] = P$ holds exactly.
        This eliminates all spurious hits from the output.
\end{enumerate}

\begin{formaldetails}[Rolling hash recurrence]
\begin{lemma}
  The rolling hash formula correctly maintains
  $t \equiv t_s \pmod{q}$ throughout the algorithm.
\end{lemma}
\begin{proof}
  After the preprocessing loop, $t = t_1 \bmod q$ by Horner's rule.
  For the inductive step, suppose $t = t_s \bmod q$ at the start of
  iteration~$s$.
  The update computes:
  \[
    t' = \bigl((t_s - T[s]\cdot d^{m-1})\cdot d + T[s+m]\bigr) \bmod q.
  \]
  But $t_{s+1} = (t_s - T[s]\cdot d^{m-1})\cdot d + T[s+m]$ by the
  definition of the window values, so $t' = t_{s+1} \bmod q$.
\end{proof}
\end{formaldetails}

\subsection{Complexity}

The running time depends on the number of spurious hits.

\paragraph{Preprocessing.}
Computing $h$, $p$, and $t_1$ takes $\Oh{m}$ time.

\paragraph{Scanning.}
In each of the $n - m + 1$ iterations, the rolling hash update and the
modular comparison take $\Oh{1}$.
The expensive part is the verification upon a hit: each verification
costs $\Oh{m}$.

Let $\nu$ denote the number of true occurrences and $c$ the number of
spurious hits.
The total running time is:
\[
  \Oh{m} + \Oh{n - m} + (\nu + c) \cdot \Oh{m}
  \;=\;
  \Oh{n + m\,(\nu + c)}.
\]

\paragraph{Worst case.}
In the worst case, \emph{every} position is a hit (true or spurious),
giving $\nu + c = n - m + 1$ and a total time of $\Oh{nm}$---no
better than the naive algorithm.
This happens, for instance, when $q = 2$, $P = \texttt{aa}$, and
$T = \texttt{aaa\ldots a}$: every window matches both modulo~$q$ and
exactly.

\paragraph{Expected case.}
The key insight is that by choosing $q$ to be a sufficiently large
\emph{random prime}, the probability of a spurious hit becomes very
small.

\begin{definition}[Fingerprint]
  For a positive integer $q$, the \textbf{fingerprint} of $P$ is
  $H_q(P) = p \bmod q$, and the fingerprint of window $s$ is
  $H_q(T_s) = t_s \bmod q$.
  A \textbf{false match} at position $s$ occurs when
  $H_q(P) = H_q(T_s)$ but $P \neq T[s, s+m-1]$.
\end{definition}

If $P \neq T[s, s+m-1]$, then $p \neq t_s$, so $q$ can only produce a
false match if $q$ divides $|p - t_s|$.
Since $|p - t_s| < d^m$, it has at most $\pi(d^m)$ distinct prime
divisors, where $\pi(u)$ denotes the number of primes up to $u$.

If we pick $q$ uniformly at random among all primes up to some bound
$I$, the probability that $q$ divides $|p - t_s|$ at any single
position~$s$ is at most $\pi(d^m) / \pi(I)$.

\begin{theorem}[Karp--Rabin, 1987]
  \label{thm:rabin-karp-error}
  If $q$ is a random prime in $\{1, \ldots, I\}$ with $I = nm^2$,
  then the probability of a false match at \emph{any single} position
  is at most $\Oh{1/m}$.
\end{theorem}

\aside{%
  By the prime number theorem,
  $\pi(u) \approx u / \ln u$, so the bound works out to roughly
  $\frac{\pi(d^m)}{\pi(nm^2)} \approx \frac{m \ln d}{\ln(nm^2)}
  \cdot \frac{1}{m} = \Oh{1/m}$ for $d$ constant.}

When $q$ is large enough (say $q \ge nm$) and chosen as a random
prime, the expected number of spurious hits across all $n - m + 1$
positions is $\Oh{n/m}$, giving an expected total time of:
\[
  \Oh{n + m} + \Oh{n/m}\cdot\Oh{m} = \Oh{n + m}.
\]
If the number of true occurrences $\nu$ is constant, this is
\textbf{expected linear time}.

\subsection{Las~Vegas vs.\ Monte~Carlo}

The Rabin--Karp algorithm as presented above always verifies hits, so it
\emph{never reports a false positive}.
This makes it a \textbf{Las~Vegas} randomised algorithm: the output is
always correct, but the running time is a random variable (it depends on
how many spurious hits occur).

\begin{description}[style=nextline, leftmargin=2cm]
  \item[Las~Vegas variant.]
    Always verify hits.
    Output is \emph{always correct}.
    Worst-case time $\Oh{nm}$, expected time $\Oh{n + m}$.

  \item[Monte~Carlo variant.]
    \emph{Skip} the verification step: when
    $t_s \bmod q = p \bmod q$, report position~$s$ as an occurrence
    without checking.
    Running time is $\Oh{n + m}$ \emph{always} (deterministic), but
    the output may contain false positives.
    The probability of any false positive is bounded by
    Theorem~\ref{thm:rabin-karp-error}.
\end{description}

\begin{remark}
  We can further reduce the error probability by using $k$ independent
  random primes $q_1, \ldots, q_k$ and declaring a match only when
  \emph{all} $k$ fingerprints agree.
  The probability of a false match drops to
  $\Oh{(1/m)^k}$---exponentially small in $k$---at the cost of
  a factor $k$ in running time.
\end{remark}

\subsection{Properties}

\begin{center}
\begin{tabular}{lcccc}
  \toprule
  & Online & Real-time & Blind & Succinct \\
  \midrule
  Rabin--Karp (Las~Vegas)
    & \checkmark & $\times$ & $\times$ & \checkmark \\
  Rabin--Karp (Monte~Carlo)
    & \checkmark & \checkmark & \checkmark & \checkmark \\
  \bottomrule
\end{tabular}
\end{center}

The algorithm is \textbf{online}: it processes $T$ left-to-right.
The Monte~Carlo variant is even \textbf{real-time} (each character
costs $\Oh{1}$ amortised) and \textbf{blind} (no backtracking), since
it never re-reads characters of $T$.
The Las~Vegas variant can backtrack during verification, so it is
neither real-time nor blind.

Both variants are \textbf{succinct}: they use only $\Oh{1}$ extra space
beyond storing $P$ and $T$ (just the variables $h$, $p$, $t$, and
the current position).

\aside{The Rabin--Karp approach generalises naturally to
\emph{multi-pattern} matching: compute fingerprints for all patterns
and store them in a hash set.
Each position of $T$ is checked against the set in $\Oh{1}$ expected
time.
This gives expected $\Oh{n + km}$ time for $k$ patterns of length $m$,
without building an Aho--Corasick automaton.}
